\section*{第 \ref{ch:b2:unicellular-diversity} 章——单细胞生物的多样性}

\begin{solution}{exo:b2:unicellular-diversity:1}
单细胞：一个细胞执行全部功能，独自生活；
群体：相同的细胞相互连接，但每个细胞都能独自生活；
多细胞：多种细胞类型，相互依存，只有生殖系
能够繁殖。\emph{E.~coli} 与酵母都是单细胞（芽体
脱落后独自生活）；\emph{Gonium} 为群体（16 个相同
细胞，各能建立一个群体）；\omterm{prop:b2:unicellular-diversity:volvox}{\emph{Volvox}} 为多细胞
（会死亡的体细胞，少数生殖细胞）。
\end{solution}

\begin{solution}{exo:b2:unicellular-diversity:2}
$S/V = 3/r$：球菌为 \qty{6}{\micro m^{-1}}，
\omterm{def:b2:unicellular-diversity:protist}{原生生物}为 \qty{0.06}{\micro m^{-1}}，相差百倍。
球菌单位体积的呼吸作用更快：其细胞质的每一部分
都在不足一微米之内靠近表面，
氧气和食物正是经由该表面进入；而\omterm{def:b2:unicellular-diversity:protist}{原生生物}的内部
经单位体积表面积少了百倍的途径获得供给。
\end{solution}

\begin{solution}{exo:b2:unicellular-diversity:3}
无细胞核（仅有\omterm{def:b2:unicellular-diversity:prokaryote}{拟核}），与具核膜相对；无线粒体或
内膜系统，呼吸链位于质膜上，
与具细胞器相对；鞭毛是由鞭毛蛋白构成的刚性螺旋，
由质子驱动旋转，与 $9+2$ 纤毛相对，后者
靠动力蛋白和 ATP 弯曲；另外还有大小（\qty{1}{\micro m}，对比数十微米）以及
\omterm{prop:b2:unicellular-diversity:gram}{肽聚糖}\omterm{def:b2:unicellular-diversity:prokaryote}{细胞壁}。细菌不能吞噬：其\omterm{def:b2:unicellular-diversity:prokaryote}{细胞壁}
不允许包裹颗粒，而变形虫借助伪足即可做到。
\end{solution}

\begin{solution}{exo:b2:unicellular-diversity:4}
纤毛：运动并将食物扫向口腔；口沟：
将颗粒汇集到胞咽；食物泡：消化摄入的
细菌，借助溶酶体酶；\omterm{prop:b2:unicellular-diversity:osmoregulation}{伸缩泡}：排出
通过渗透作用进入的水；大核：工作核，
数百份基因拷贝被转录，用于日常生活；小核：
生殖二倍体核，用于减数分裂和接合。
\end{solution}

\begin{solution}{exo:b2:unicellular-diversity:5}
$t = x^2/2D$：细菌中为 $(10^{-6})^2/10^{-9} = \qty{1}{ms}$；
变形虫中为 $(5\times 10^{-4})^2/10^{-9} = \qty{250}{s}$，约四分钟。
变形虫不能仅依靠扩散来输送代谢物：
它需要胞质环流以及马达驱动的
沿细胞骨架的运输，而细菌无需这些。
\end{solution}

\begin{solution}{exo:b2:unicellular-diversity:6}
$Re = \rho vL/\eta$：\emph{Paramecium} 为 $10^3\times 10^{-3}\times
2\times 10^{-4}/10^{-3} = 0.2$；细菌为 $10^3\times 2\times
10^{-5}\times 2\times 10^{-6}/10^{-3} = 4\times 10^{-5}$；游动生物为
$10^3\times 1\times 2/10^{-3} = 2\times 10^6$。当 $Re \gg 1$ 时，惯性
使游动生物在划水后仍继续前进；当 $Re \ll 1$ 时，黏性阻力
使细菌在一纳米之内停止，因此它只在
摆动时才运动。
\end{solution}

\begin{solution}{exo:b2:unicellular-diversity:7}
$v_s = 2r^2\Delta\rho g/9\eta = 2\times 10^{-10}\times 100\times
9.81/(9\times 10^{-3}) = \qty{2.2e-5}{m/s}$；下沉 \qty{50}{m} 需要
$50/2.2\times 10^{-5} = \qty{2.3e6}{s}$，即 27 天。半径减半
使 $v_s$ 降为原来的四分之一，需 108 天。将 $\Delta\rho$ 降至
\qty{20}{kg/m^3}，使其降为原来的五分之一，需 135 天。
\end{solution}

\begin{solution}{exo:b2:unicellular-diversity:8}
$V = \tfrac{4}{3}\pi\times 100\times 25\times 25 =
\qty{2.6e5}{\micro m^3}$。在 \qty{1800}{s} 内：$2.6\times 10^5/1800 =
\qty{145}{\micro m^3/s}$；因为 $\qty{1}{\micro m^3} =
\qty{1e-15}{L}$，即 \qty{1.5e-13}{L/s}。
\end{solution}

\begin{solution}{exo:b2:unicellular-diversity:9}
在 \qty{1}{mm} 内加倍，相当于每
微米增加约 \qty{0.07}{\%}（因 $2^{1/1000} - 1 = 0.0007$），
因此跨越整个细胞为 \qty{0.07}{\%}。一次 \qty{1}{s} 的游动覆盖 \qty{20}{\micro m}：
\qty{1.4}{\%}。要分辨 \qty{0.07}{\%} 的差异，细胞需要
在两端各计数约 $(1/0.0007)^2 = 2\times 10^6$ 个分子，
远多于其数千个受体一秒内能结合的分子数；
\qty{1.4}{\%} 只需约 \num{5000} 个，这是可行的。时间上的比较
把无法实现的空间测量转化为可行的
测量，方法是让游动过程完成积分。
\end{solution}

\begin{solution}{exo:b2:unicellular-diversity:10}
表面积 $4\pi r^2$，其中 $r = \qty{250}{\micro m}$：\qty{7.9e5}{\micro
m^2}；细胞质体积 $\approx$ 表面积 $\times$ \qty{2}{\micro m}
$= \qty{1.6e6}{\micro m^3}$；比值为 \qty{0.5}{\micro m^{-1}}。对于
\emph{E.~coli}，视为圆柱体：$S/V = 2\pi r L/\pi r^2 L = 2/r =
\qty{4}{\micro m^{-1}}$。该巨型细菌仅比
\emph{E.~coli} 差八倍，而非 500 倍，因为其细胞质的每一点
都在距表面 \qty{2}{\micro m} 之内。该液泡储备
硝酸盐（其电子受体），使其能
在两次带来硝酸盐的罕见搅动之间，于缺氧沉积物中以硫化物呼吸
（\cref{ch:b2:microbial-metabolism}）。
\end{solution}

\begin{solution}{exo:b2:unicellular-diversity:11}
细胞不能同时游动和分裂。由 16 个细胞组成的群体，若在分裂所需的数小时内停止
游动，则按斯托克斯定律下沉的
距离可忽略（其半径为数十微米，
下沉速度为每秒微米量级）；整个群体可以分裂
而不会有任何损失。一个由 \num{10000} 个细胞组成、半径为
\qty{500}{\micro m} 的球体下沉速度会快百倍——
每秒数毫米、每小时数米——移出光照区：
因此让多数细胞保持游动、少数分裂更为有利。生殖细胞
之所以大，是因为它们必须携带资源以构建整个
由数千个细胞组成的子代球体，经连续分裂
而不摄食，因此由体细胞为其提供补给。
\end{solution}

\begin{solution}{exo:b2:unicellular-diversity:12}
演化支由一个祖先及其全部后代组成；等级则是若干谱系各自达到的
组织层次。\omterm{def:b2:unicellular-diversity:protist}{原生生物}——绿藻（与陆生植物同属）、硅藻（与褐藻同属）、
纤毛虫（与甲藻及 \emph{Plasmodium} 同属）、变形虫（与动物和真菌相近）——
它们只共享独居的真核细胞这一特征，
其最近共同祖先即全部真核生物（包括我们）的祖先，
因此是一个等级。
“\omterm{def:b2:unicellular-diversity:prokaryote}{原核生物}”以其所缺乏的特征联结细菌与古菌；古菌与真核生物的
亲缘关系比与细菌更近，因此这一术语同样是一个等级。
蓝细菌的后代均源自一个发明了产氧光合作用的祖先，
并包括其全部后代（叶绿体谱系除外）：
这是一个演化支。
\end{solution}

\begin{solution}{pb:b2:unicellular-diversity:1}
\textbf{1.} $1\times 1\times 0.1 = \qty{0.1}{mm^3} = \qty{0.1}{\micro L}$。
\textbf{2.} $24/\qty{0.1}{\micro L} = 240$ 每 \unit{\micro L} $=
\qty{2.4e5}{}$ 个细胞每毫升。
\textbf{3.} $V = \tfrac{4}{3}\pi\times 125 = \qty{524}{\micro m^3}$；
$S/V = 3/5 = \qty{0.6}{\micro m^{-1}}$。
\textbf{4.} $2.4\times 10^5\times 524 = \qty{1.26e8}{\micro m^3}$ 每
毫升，而 $\qty{1}{mL} = \qty{1e12}{\micro m^3}$：占比为
$1.3\times 10^{-4}$，即 \qty{0.013}{\%}。
\textbf{5.} $\qty{1000}{m^3} = \qty{1e9}{mL}$：$2.4\times 10^{14}$
个细胞。
\textbf{6.} 蓝细菌是由小得多的细胞（约几微米）组成的丝状体，
呈蓝绿色，没有明显的叶绿体或
细胞核，且不借助鞭毛游动；藻类则为单个
卵形细胞，具两根鞭毛、一个眼点和一个杯状
叶绿体。
\textbf{7.} $384/24 = 16 = 2^4$：\qty{48}{h} 内四次倍增，
$T = \qty{12}{h}$。
\textbf{8.} $k = \ln 2/T = 0.693/12 = \qty{0.058}{h^{-1}}$。
\textbf{9.} $t = \ln(10^8/2.4\times 10^5)/k = \ln(417)/0.058 =
6.03/0.058 = \qty{104}{h}$，约 4.3 天。
\textbf{10.} 营养物质（磷酸盐、硝酸盐）耗尽；细胞相互遮荫，
使每个细胞所获光照减少；摄食者（纤毛虫、
轮虫、\emph{Daphnia}）在其上大量繁殖；细胞数达到每毫升 \qty{1e8}{}
个时，细胞将占据约百分之五的体积，
并耗尽溶解的 $\mathrm{CO_2}$。
\textbf{11.} 仅有一半时间在生长：平均速率减半，
为 \qty{208}{h}，即 8.7 天。
\textbf{12.} 计数反映的是种群，而非单个细胞：
若一个细胞每 \qty{36}{h} 释放 8 个子代，则种群
在三个各为 \qty{12}{h} 的\omterm{def:b2:unicellular-diversity:division}{世代时间}内增加为原来的 8 $= 2^3$ 倍。
\omterm{def:b2:unicellular-diversity:division}{世代时间}即数量的平均倍增时间，
无论分裂以何种方式分组。
\textbf{13.} $Re = 10^3\times 10^{-4}\times 10^{-5}/10^{-3} = 10^{-3}$。
\textbf{14.} $m = 1050\times 5.24\times 10^{-16} = \qty{5.5e-13}{kg}$；
$6\pi\eta r = 6\pi\times 10^{-3}\times 5\times 10^{-6} =
\qty{9.4e-8}{kg/s}$；$d = 5.5\times 10^{-13}\times 10^{-4}/9.4\times
10^{-8} = \qty{5.8e-10}{m}$：约半纳米。
\textbf{15.} $v_s = 2\times 25\times 10^{-12}\times 50\times
9.81/(9\times 10^{-3}) = \qty{2.7e-6}{m/s}$，\qty{2.7}{\micro m/s}。
\textbf{16.} 下沉 \qty{1}{m}：$1/2.7\times 10^{-6} =
\qty{3.7e5}{s}$，4.3 天；以 \qty{100}{\micro m/s} 的速度向上游动：
\qty{1e4}{s}，2.8 小时。游动快 37 倍。
\textbf{17.} $c_\infty = \qty{1e-2}{mol/m^3}$：$J = 4\pi\times
1.9\times 10^{-9}\times 5\times 10^{-6}\times 10^{-2} =
\qty{1.2e-15}{mol/s} = 7.2\times 10^8$ 个分子每秒。
\textbf{18.} \qty{100}{pg} $= \qty{1e-10}{g} = \qty{8.3e-12}{mol}$ 的碳，
在 \qty{12}{h} $= \qty{43200}{s}$ 内翻倍：
\qty{1.9e-16}{mol/s} $= 1.2\times 10^8$ 个原子每秒。供给与
需求之比：$7.2/1.2 = 6$。
\textbf{19.} 供给增至十倍（$\propto r$），
需求增至一千倍（$\propto r^3$）：比值变为 $6/100 = 0.06$；
这样的细胞仅能以原速率的百分之六生长，因此，
这种大小、依靠扩散供给的藻类
若无搅拌、内部运输或代谢显著减慢，便不可能存在。
\textbf{20.} $\Delta\Pi = RT\Delta c = 8.314\times 293\times 99 =
\qty{2.4e5}{Pa}$，\qty{2.4}{bar}。
\textbf{21.} 面积 $4\pi(5\times 10^{-6})^2 = \qty{3.14e-10}{m^2}$；
流入量 $L_p A\Delta\Pi = 10^{-14}\times 3.14\times 10^{-10}\times
2.4\times 10^5 = \qty{7.6e-19}{m^3/s} = \qty{0.76}{\micro m^3/s}$。
\textbf{22.} $524/0.76 = \qty{690}{s}$，约十一分钟。
\textbf{23.} $P = \Delta\Pi\times$ 流量 $= 2.4\times 10^5\times
7.6\times 10^{-19} = \qty{1.8e-13}{W}$。
\textbf{24.} 一个 ATP：$5\times 10^4/6.02\times 10^{23} =
\qty{8.3e-20}{J}$：$1.8\times 10^{-13}/8.3\times 10^{-20} = 2.2\times
10^6$ 个 ATP 每秒。碳固定：$3\times 1.2\times 10^8 =
3.6\times 10^8$ 个 ATP 每秒。液泡耗能约占
其中的 \qty{0.6}{\%}——以较小代价防止细胞胀裂。
\textbf{25.} 水的流入量为 \qty{0.76}{\micro m^3/s}（即细胞体积在
十一分钟内的变化量）；若无\omterm{prop:b2:unicellular-diversity:osmoregulation}{伸缩泡}，细胞将
约在 \qty{700}{s} 内体积翻倍；维持体内干燥至少需要
每秒 $2\times 10^6$ 个 ATP，不足其
光合作用固定碳所耗能量的百分之一。
\end{solution}
